\section{Limitations}
\label{sec:limitations}

This is a deliberately narrow, honest study and should be read as such.
\begin{itemize}
  \item \textbf{Scale of models and data.} The main study uses \texttt{qwen3:8b}
  with \texttt{nomic-embed}/\texttt{MiniLM} on LoCoMo. \S\ref{sec:generality}
  extends the two load-bearing findings to a second judge (\texttt{llama3.2:3b}),
  a second answerer (\texttt{llama3.2:3b}), and a second dataset
  (LongMemEval~\cite{wu2024longmemeval}), but all models are small and local;
  whether the conclusions hold for frontier models is untested. Absolute numbers
  are not competitive with frontier systems and are not meant to be---only the
  \emph{relative, within-pipeline} comparisons are claimed.
  \item \textbf{LLM-as-judge.} Scoring uses the same model family as generation;
  self-preference bias is possible. We mitigate by comparing arms under an
  identical judge, but a human-validated subset is future work.
  \item \textbf{Metric thresholds.} \mdr/\rtf{} are threshold- and
  embedder-dependent; we report curves rather than single numbers, and do not
  compare absolute \mdr{} across embedders.
  \item \textbf{Deduplication scope.} We test one blind-merge threshold and one
  redundancy-aware composer (MMR). A learned or store-level query-distribution
  method could differ, though our composition-level negative gives little reason
  to expect a large gain.
  \item \textbf{Faithfulness probe.} The $1/400$ result is a single-judge,
  single-extractor, conservative-parse probe; it bounds the effect as small in
  \emph{this} grounded setup, not in generative or tool-use memory.
  \item \textbf{No new method.} We claim a diagnosis, not a system that improves
  the state of the art.
\end{itemize}
